\section{Input, quá trình và output của sáu lựa chọn nội suy ngược}
\subsection*{Input: các ô đúng tên trên giao diện}
Bảng này được lấy từ biểu mẫu của tool. Giá trị mặc định là ví dụ nhập, cần thay bằng dữ liệu đề. Hai dãy $x,y$ là bảng dữ liệu riêng bên dưới nếu chức năng dùng bảng.
\subsubsection*{Ô chung của các lựa chọn trong nhóm}
\begin{center}\small\begin{tabular}{@{}p{.32\linewidth}p{.17\linewidth}p{.43\linewidth}@{}}\toprule
Ô nhập & Khóa JSON & Cách nhập / mặc định\\\midrule
Mốc dùng (JSON,\allowbreak{} trống: toàn bộ bảng) & \texttt{selected\_indices} & Để trống
\par Ví dụ [9,\allowbreak{}10,\allowbreak{}11,\allowbreak{}12] dùng đúng bốn dòng đó,\allowbreak{} đánh số từ 0;\allowbreak{} dữ liệu nguồn vẫn được giữ.\\[3pt]
Mức trình bày & \texttt{detail\_level} & \texttt{Tổng hợp}
\par Bảng tổng hợp luôn hiện;\allowbreak{} có thể mở từng nhóm phép tính hoặc từng vòng. Bản xuất vẫn giữ đủ chi tiết.\\[3pt]
\bottomrule\end{tabular}\end{center}
\subsubsection*{Ô riêng: Lặp tiến}
\begin{center}\small\begin{tabular}{@{}p{.32\linewidth}p{.17\linewidth}p{.43\linewidth}@{}}\toprule
Ô nhập & Khóa JSON & Cách nhập / mặc định\\\midrule
Giá trị y cần tìm ngược & \texttt{target} & \texttt{4}
\par Tập cặp nhập là các mốc đã chọn;\allowbreak{} bậc bằng số cặp trừ 1. Không tự đổi nhánh hoặc mốc. Lặp thực hành mặc định kiểm tra thêm phần dư 1e-6,\allowbreak{} tối đa 100 vòng.\\[3pt]
Cận x trái (trống: khoảng đổi dấu duy nhất) & \texttt{a} & Để trống\\[3pt]
Cận x phải (trống: khoảng đổi dấu duy nhất) & \texttt{b} & Để trống\\[3pt]
Dung sai vị trí/\allowbreak{}bước theo x & \texttt{tol} & \texttt{1e-06}
\par Giữ toàn bộ tập mốc và gốc đầu/\allowbreak{}cuối tập;\allowbreak{} khoảng nghiệm có thể ở giữa bảng.\\[3pt]
Dung sai phần dư & \texttt{residual\_tol} & \texttt{1e-06}\\[3pt]
Số vòng tối đa & \texttt{max\_iter} & \texttt{100}\\[3pt]
Quy tắc khởi tạo & \texttt{initial\_rule} & \texttt{Tuyến tính tại gốc}\\[3pt]
t₀ khi khởi tạo thủ công & \texttt{initial\_t} & Để trống\\[3pt]
\bottomrule\end{tabular}\end{center}
\subsubsection*{Ô riêng: Lặp lùi}
\begin{center}\small\begin{tabular}{@{}p{.32\linewidth}p{.17\linewidth}p{.43\linewidth}@{}}\toprule
Ô nhập & Khóa JSON & Cách nhập / mặc định\\\midrule
Giá trị y cần tìm ngược & \texttt{target} & \texttt{2}
\par Tập cặp nhập là các mốc đã chọn;\allowbreak{} bậc bằng số cặp trừ 1. Không tự đổi nhánh hoặc mốc. Lặp thực hành mặc định kiểm tra thêm phần dư 1e-6,\allowbreak{} tối đa 100 vòng.\\[3pt]
Cận x trái (trống: khoảng đổi dấu duy nhất) & \texttt{a} & Để trống\\[3pt]
Cận x phải (trống: khoảng đổi dấu duy nhất) & \texttt{b} & Để trống\\[3pt]
Dung sai vị trí/\allowbreak{}bước theo x & \texttt{tol} & \texttt{1e-06}
\par Giữ toàn bộ tập mốc và gốc đầu/\allowbreak{}cuối tập;\allowbreak{} khoảng nghiệm có thể ở giữa bảng.\\[3pt]
Dung sai phần dư & \texttt{residual\_tol} & \texttt{1e-06}\\[3pt]
Số vòng tối đa & \texttt{max\_iter} & \texttt{100}\\[3pt]
Quy tắc khởi tạo & \texttt{initial\_rule} & \texttt{Tuyến tính tại gốc}\\[3pt]
t₀ khi khởi tạo thủ công & \texttt{initial\_t} & Để trống\\[3pt]
\bottomrule\end{tabular}\end{center}
\subsubsection*{Ô riêng: Lặp có cận co}
\begin{center}\small\begin{tabular}{@{}p{.32\linewidth}p{.17\linewidth}p{.43\linewidth}@{}}\toprule
Ô nhập & Khóa JSON & Cách nhập / mặc định\\\midrule
Giá trị y cần tìm ngược & \texttt{target} & \texttt{4}
\par Tập cặp nhập là các mốc đã chọn;\allowbreak{} bậc bằng số cặp trừ 1. Không tự đổi nhánh hoặc mốc. Lặp thực hành mặc định kiểm tra thêm phần dư 1e-6,\allowbreak{} tối đa 100 vòng.\\[3pt]
Cận x trái của khoảng nghiệm & \texttt{a} & \texttt{1.6}\\[3pt]
Cận x phải của khoảng nghiệm & \texttt{b} & \texttt{1.8}\\[3pt]
Dung sai vị trí & \texttt{tol} & \texttt{1e-06}\\[3pt]
Dung sai phần dư & \texttt{residual\_tol} & \texttt{1e-06}\\[3pt]
Hướng nội suy & \texttt{direction} & \texttt{Tiến}\\[3pt]
Số vòng tối đa & \texttt{max\_iter} & \texttt{100}\\[3pt]
Quy tắc khởi tạo & \texttt{initial\_rule} & \texttt{Tuyến tính tại gốc}\\[3pt]
t₀ khi khởi tạo thủ công & \texttt{initial\_t} & Để trống\\[3pt]
\bottomrule\end{tabular}\end{center}
\subsubsection*{Ô riêng: Hàm ngược — Newton /\allowbreak{} Hàm ngược — Lagrange}
\begin{center}\small\begin{tabular}{@{}p{.32\linewidth}p{.17\linewidth}p{.43\linewidth}@{}}\toprule
Ô nhập & Khóa JSON & Cách nhập / mặc định\\\midrule
Giá trị y cần tìm ngược & \texttt{target} & \texttt{6}
\par Tập cặp nhập là các mốc đã chọn;\allowbreak{} bậc bằng số cặp trừ 1. Không tự đổi nhánh hoặc mốc. Lặp thực hành mặc định kiểm tra thêm phần dư 1e-6,\allowbreak{} tối đa 100 vòng.\\[3pt]
Số mốc đã chọn (để trống: toàn tập nhập) & \texttt{count} & Để trống
\par Không tự lấy mốc gần mức. Muốn dùng ít mốc hơn,\allowbreak{} giữ đúng các cặp cần dùng trong bảng nhập.\\[3pt]
\bottomrule\end{tabular}\end{center}
\subsubsection*{Ô riêng: Chia đôi đa thức}
\begin{center}\small\begin{tabular}{@{}p{.32\linewidth}p{.17\linewidth}p{.43\linewidth}@{}}\toprule
Ô nhập & Khóa JSON & Cách nhập / mặc định\\\midrule
Giá trị y cần tìm ngược & \texttt{target} & \texttt{2}
\par Tập cặp nhập là các mốc đã chọn;\allowbreak{} bậc bằng số cặp trừ 1. Không tự đổi nhánh hoặc mốc. Lặp thực hành mặc định kiểm tra thêm phần dư 1e-6,\allowbreak{} tối đa 100 vòng.\\[3pt]
Cận trái & \texttt{a} & \texttt{1}\\[3pt]
Cận phải & \texttt{b} & \texttt{2}\\[3pt]
Dung sai vị trí & \texttt{tol} & \texttt{0.001}\\[3pt]
\bottomrule\end{tabular}\end{center}

\paragraph{Input chung.} $x,y$ là bảng nguồn; $\bar y$ là mức cần tìm $x$. Ô \emph{Mốc dùng} chọn chỉ số từ 0, trống dùng toàn bảng. Khoảng $[a,b]$ là nơi muốn tìm nghiệm, không thay gốc của đa thức. Không lấy số mốc làm số vòng lặp.

\subsection{Bước đầu chung: bảng khảo sát tổng quát}
\begin{center}\small\begin{tabular}{cllll}\toprule
Dòng&$x_i$&$F_i=y_i-\bar y$&So với dòng sau&Thông tin chuyển tiếp\\\midrule
0&$x_0$&$y_0-\bar y$&Dấu $F_0,F_1$&Mốc đúng mức/khoảng đổi dấu\\
$i$&$x_i$&$y_i-\bar y$&Dấu $F_i,F_{i+1}$&Ứng viên $[x_i,x_{i+1}]$\\
$n$&$x_n$&$y_n-\bar y$&Không có dòng sau&Kiểm riêng $F_n=0$\\\bottomrule
\end{tabular}\end{center}
Sau khảo sát mới chọn nhánh, tập mốc và thuật toán. Nhiều khoảng thì phải giải từng nhánh cần tìm. Bảng chỉ phát hiện đổi dấu giữa các mẫu; có thể bỏ sót nghiệm tiếp xúc.

\subsection{Hàm ngược -- Newton và Hàm ngược -- Lagrange}
\textbf{Input}: tập mốc thuộc nhánh tăng/giảm nghiêm ngặt, tung độ phân biệt; mức trong dải tung độ nếu không bật ngoại suy. \textbf{Quá trình}: đảo cặp $(u_i,v_i)=(y_i,x_i)$ $\to$ dựng $Q(y)$ $\to$ thế $\bar y$.
\begin{center}\small\begin{tabular}{cllll}\toprule
Dòng&$u_i$&$v_i$&Tỉ sai phân cấp 1&Cấp 2\\\midrule
0&$y_0$&$x_0$&$(x_1-x_0)/(y_1-y_0)$&$D_{0,2}$\\
$i$&$y_i$&$x_i$&$(x_{i+1}-x_i)/(y_{i+1}-y_i)$&$D_{i,2}$ nếu đủ mốc\\
$n$&$y_n$&$x_n$&&\\\bottomrule
\end{tabular}\end{center}
Newton lấy $D_{0,k}$ rồi cộng
\[
Q(\bar y)=x_0+D_{0,1}(\bar y-y_0)+D_{0,2}(\bar y-y_0)(\bar y-y_1)+\cdots.
\]
Lagrange tính $\ell_i(\bar y)=\prod_{j\ne i}(\bar y-y_j)/(y_i-y_j)$ rồi cộng
\[
\widehat x=x_0\ell_0(\bar y)+x_1\ell_1(\bar y)+\cdots+x_n\ell_n(\bar y).
\]
\textbf{Output}: đa thức hàm ngược $Q$, giá trị $Q(\bar y)$ và tập mốc; không có vòng hội tụ. Tung độ thường không đều dù hoành độ gốc đều.

\subsection{Lặp tiến, Lặp lùi, Lặp có cận co}
\textbf{Input}: lưới đều, tập mốc cố định, mức, khoảng, ngưỡng vị trí $\varepsilon_x$, ngưỡng phần dư $\varepsilon_y$, khởi tạo và số vòng tối đa. Lặp có cận co thêm lựa chọn hướng; hai mục tiến/lùi mặc định dùng tiêu chuẩn thực hành.

\textbf{Bước 1}: dựng đa thức thuận $P_t(t)=c_0+dt+H(t)$ một lần:
\begin{align*}
\text{Tiến: }&t=(x-x_0)/h,\ c_0=y_0,\ d=\Delta y_0,\\
&H(t)=\tfrac{\Delta^2y_0}2t(t-1)+\tfrac{\Delta^3y_0}6t(t-1)(t-2)+\cdots;\\
\text{Lùi: }&t=(x-x_n)/h,\ c_0=y_n,\ d=\Delta y_{n-1},\\
&H(t)=\tfrac{\Delta^2y_{n-2}}2t(t+1)+\tfrac{\Delta^3y_{n-3}}6t(t+1)(t+2)+\cdots.
\end{align*}
\textbf{Bước 2}: lập $\varphi(t)=(\bar y-c_0-H(t))/d$, chọn $t_0$ theo ô khởi tạo. $d$ phải khác 0. Nếu có chứng nhận, kiểm $t_0\in J$, $\varphi(J)\subseteq J$, $q\ge\sup_J|\varphi'|$, $q<1$.

\textbf{Bước 3}: điền bảng vòng lặp tổng quát theo hai nửa để đọc rõ:
\begin{center}\small\begin{tabular}{cllll}\toprule
Vòng&Giá trị nhận&Phần bậc cao&Tử số&Giá trị mới\\\midrule
1&$t_0$&$H(t_0)$&$\bar y-c_0-H(t_0)$&$t_1=(\bar y-c_0-H(t_0))/d$\\
$k+1$&$t_k$&$H(t_k)$&$\bar y-c_0-H(t_k)$&$t_{k+1}=(\bar y-c_0-H(t_k))/d$\\\bottomrule
\end{tabular}\end{center}
\begin{center}\small\begin{tabular}{cllll}\toprule
Vòng&Đổi lại hoành độ&Bước hoặc cận&Phần dư tại điểm mới&Quyết định\\\midrule
$k+1$&$x_{k+1}=a_0+ht_{k+1}$&$D_{k+1}$ hoặc $B_{k+1}$&$|P_t(t_{k+1})-\bar y|$&Đạt cả hai ngưỡng?\\\bottomrule
\end{tabular}\end{center}
Ở đây $a_0$ là gốc đầu/cuối tập mốc, không phải cận trái khoảng nghiệm.
$D_{k+1}=|h||t_{k+1}-t_k|$ dùng ở thực hành;
$B_{k+1}=|h|q|t_{k+1}-t_k|/(1-q)$ dùng khi có chứng nhận hợp lệ.
Vòng chưa đạt chuyển \emph{giá trị mới} thành đầu vào vòng sau. Không dùng $H(t_k)$ để tính phần dư tại $t_{k+1}$.

\textbf{Output}: nghiệm mô hình trong khoảng yêu cầu, trạng thái dừng, gốc/hướng/tập mốc, bảng mọi vòng, bước/cận và phần dư. Hết vòng hoặc tràn số là trạng thái chưa đạt. Đúng một mốc hoặc đa thức thực tế tuyến tính có thể cho kết quả không cần vòng lặp.

\subsection{Chia đôi đa thức}
\textbf{Input}: cùng bảng và mức nhưng cần khoảng $[a,b]$, ngưỡng vị trí. \textbf{Quá trình}: dựng $P$ $\to$ đặt $F=P-\bar y$ $\to$ kiểm hai đầu $\to$ thu khoảng:
\begin{center}\small\begin{tabular}{cllll}\toprule
Lần&Khoảng nhận&Trung điểm&Giá trị kiểm&Khoảng chuyển tiếp\\\midrule
$k+1$&$[a_k,b_k]$&$c_k=(a_k+b_k)/2$&$F(c_k)$&Giữ nửa có đổi dấu\\\bottomrule
\end{tabular}\end{center}
Khi nửa độ rộng đạt ngưỡng, \textbf{output} là trung điểm của khoảng cuối, khoảng cuối và cận nửa độ rộng. Trung điểm vừa thử ở hàng cuối không nhất thiết là trung điểm của khoảng sau thu. Phần giải thích và ví dụ đầy đủ theo từng mô hình ở các mục tiếp theo.
